% -*- TeX-master: "../main.tex" -*-
\chapter{Analiza konkurencji}

Przed rozpoczęciem opracowywania projektu naszej aplikacji, szczegółowo przeanalizowaliśmy rynek. Pobraliśmy i przetestowaliśmy programy konkurencji, oraz przeanalizowaliśmy opinie ich użytkowników. Zebrane informacje skompilowane są poniżej.

\section{Code App}
\MyFigure[0.3\textwidth]{images/konkurencja/codeapp.png}{Aplikacja Code App}{fig:codeapp}{App Store}
\begin{itemize}
    \item \textbf{Docelowa platforma:} iOS/IPadOS
    \item \textbf{Docelowa grupa:} Programiści 
    \item \textbf{Główne przewagi konkurencyjne:} Silnik Monaco (z VS Code). Posiada wbudowane lokalne środowiska (Node.js, Python, PHP, C++) oraz klienta Git. 
    \item \textbf{Słabe strony:} Ograniczenia wynikające z Sandbox (system blokuje uruchamianie zewnętrznych procesów, systemowych usług w tle i narzędzi do wirtualizacji), brak autouzupełniania nazw klas i importów dla C++ i Java, brak dokumentacji. 
    \item \textbf{Model biznesowy:} Jednorazowa opłata w wysokości 30 zł
\end{itemize}

\newpage
\section{AndroidIDE}
\MyFigure[0.3\textwidth]{images/konkurencja/androidIDE.png}{Aplikacja AndroidIDE}{fig:androidide}{Google Play}
\begin{itemize}
    
    \item \textbf{Docelowa platforma:} Android
    \item \textbf{Docelowa grupa:} Początkujący i zaawansowani programiści aplikacji na Androida 
    \item \textbf{Główne przewagi konkurencyjne:} Pełne wsparcie dla systemu Gradle, aplikacja open source, autouzupełnianie kodu Java/Kotlin, kompilacja, budowanie i uruchamianie aplikacji w pełni lokalnie, zintegrowany projektant UI. 
    \item \textbf{Słabe strony:} Wysokie wymagania sprzętowe, projekt został porzucony w 2024. 
    \item \textbf{Model biznesowy:} Projekt open source (Bezpłatna) 
\end{itemize}

\newpage
\section{Acode}
\MyFigure[0.3\textwidth]{images/konkurencja/acode.png}{Aplikacja Acode}{fig:acode}{Google Play}
\begin{itemize}
    
    \item \textbf{Docelowa platforma:} Android 
    \item \textbf{Docelowa grupa:} Web-deweloperzy (Front-end), programiści 
    \item \textbf{Główne przewagi konkurencyjne:} Możliwość otwierania dużych plików, autouzupełnianie i kolorowanie składni dla ponad 100 języków programowania, wbudowany podgląd stron internetowych, interaktywna konsola JavaScript, wiele open source wtyczek do pobrania w aplikacji, skróty klawiszowe. 
    \item \textbf{Słabe strony:} Mało intuicyjny design (ciężko się odnaleźć w aplikacji)
    \item \textbf{Model biznesowy:} Dostępna jest wersja bezpłatna z reklamami oraz wersja płatna bez reklam za 13,99 zł. 
\end{itemize}

\newpage
\section{Textastic}
\MyFigure[0.3\textwidth]{images/konkurencja/textastic.png}{Aplikacja Textastic}{fig:textastic}{App Store}
\begin{itemize}
    \item \textbf{Docelowa platforma:} iOS/IPadOS 
    \item \textbf{Docelowa grupa:} Web-developerzy 
    \item \textbf{Główne przewagi konkurencyjne:} Świetna obsługa FTP, SFTP, WebDAV, integracja z chmurą (Dropbox, Google Drive), wbudowany terminal SSH, możliwość podglądu stron internetowych, kolorowanie składni dla ponad 80 języków, uzupełnianie kodu dla HTML, CSS, JS, C i PHP.
    \item \textbf{Słabe strony:} Brak kompilatorów, brak integracji z Git, ograniczone autouzupełnianie (uzupełnianie w oparciu o słowa, a nie o semantykę jak np. w VSC) 
    \item \textbf{Model biznesowy:} Płatna: jednorazowy zakup za 399,99 zł lub comiesięczna subskrypcja za 14,99 zł. 
\end{itemize}

\newpage
\section{Spck Editor}
\MyFigure[0.3\textwidth]{images/konkurencja/spck.png}{Aplikacja Spck Editor}{fig:spck}{Google Play}
\begin{itemize}
    
    \item \textbf{Docelowa platforma:} Android / iOS
    \item \textbf{Docelowa grupa:} Front-end dev (HTML, CSS, JS) 
    \item \textbf{Główne przewagi konkurencyjne:} Integracja z Git, asystent AI, Silnik Monaco Editor, wbudowana przeglądarka z narzędziami deweloperskimi, Spck Labs oferujące udostępnione projekty innych użytkowników do nauki i zabawy. 
    \item \textbf{Słabe strony:} Asystent AI oraz Predictive keyboard dostępne tylko dla użytkowników z miesięczną subskrypcją. 
    \item \textbf{Model biznesowy:} Bezpłatna, oraz wersje płatne: Spck Lite w cenie 35,99 zł jednorazowo (dostajemy w niej ekskluzywny motyw, dostęp do Predictive keyboard, customowe Snippety i paleta poleceń), Spck for NodeJS jako subskrypcja za 4,49 zł/miesiąc (integracja z NodeJS), subskrypcja „Supporter” za 8,79 zł/miesiąc (całkowite usunięcie reklam, dostęp do wszystkich płatnych motywów wizualnych interfejsu oraz dodatkowych opcji personalizacji środowiska pracy) oraz subskrypcja „Gold” za 23,99 zł/miesiąc (zawiera wszystkie korzyści z wersji Supporter, a dodatkowo zapewnia stały przydział waluty „Gold” przeznaczonej na funkcje premium, takie jak inteligentne autouzupełnianie kodu wspierane przez sztuczną inteligencję oraz rozszerzone możliwości synchronizacji projektów). 
\end{itemize}

\newpage
\section{Pydroid 3} 
\MyFigure[0.3\textwidth]{images/konkurencja/pydroid.png}{Aplikacja Pydroid 3}{fig:pydroid}{Google Play}
\begin{itemize}
    
    \item \textbf{Docelowa platforma:} Android
    \item \textbf{Docelowa grupa:} Osoby uczące się języka Python, Entuzjaści Data Science i AI. 
    \item \textbf{Główne przewagi konkurencyjne:} W pełni gotowy do użytku po zainstalowaniu aplikacji, posiada własne repozytorium prekompilowanych bibliotek (pozwala na bezproblemową instalację ciężkich paczek takich jak NumPy, Pandas, SciPy, Matplotlib, a także na uruchomienie interaktywnych notatników Jupyter), Obsługa interfejsów graficznych (np. Tkinter), wbudowany kompilator C/C++ (dla zaawansowanych bibliotek Pythona) 
    \item \textbf{Słabe strony:} Agresywne reklamy, zaawansowane biblioteki (np. TensorFlow) wymaga wersji Premium aplikacji, uzupełnianie i sprawdzanie kodu tylko w wersji Premium. 
    \item \textbf{Model biznesowy:} Bezpłatna, oraz wersja Premium z podziałem na subskrypcję 29,99 zł/miesiąc lub jednorazowy zakup za 99,99 zł. 
\end{itemize}

\newpage
\section{Pyto}
\MyFigure[0.3\textwidth]{images/konkurencja/pyto.png}{Aplikacja Pyto}{fig:pyto}{App Store}
\begin{itemize}
    
    \item \textbf{Docelowa platforma:} iOS 
    \item \textbf{Docelowa grupa:} Programiści Pythona w ekosystemie Apple 
    \item \textbf{Główne przewagi konkurencyjne:} Pozwala na uruchamianie skryptów prosto z aplikacji Skróty, pozwala na kodowanie własnych widżetów na ekran główny, dostęp do ciężkich bibliotek (NumPy, Pandas, SciPy, Matplotlib itd.) od razu po pobraniu aplikacji, pozwala na tworzenie interfejsów systemu iOS. 
    \item \textbf{Słabe strony:} Niestabilność aplikacji, brak darmowej wersji, ubogi interfejs UI. 
    \item \textbf{Model biznesowy:} Wersja Lite za 29,99 zł oraz Wersja Pełna za 59,99 zł. 
\end{itemize}

\newpage
\section{Carnets - Jupyter}
\MyFigure[0.3\textwidth]{images/konkurencja/cartens.png}{Aplikacja Carnets - Jupyter}{fig:carnet}{App Store}
\begin{itemize}
    
    \item \textbf{Docelowa platforma:} iOS/IPadOS
    \item \textbf{Docelowa grupa:} Analitycy danych i badacze 
    \item \textbf{Główne przewagi konkurencyjne:} Jupyter Notebook w pełni lokalnie, potężne biblioteki analityczne (NumPy, Pandas, Matplotlib, Bokeh, Seaborn) od razu zainstalowane, możliwość edycji i zapisu notatników w chmurze. 
    \item \textbf{Słabe strony:} Potrzeba zainstalowania osobnej aplikacji zawierającej bibliotekę do uczenia maszynowego (Scipy), przestarzały wygląd UI
    \item \textbf{Model biznesowy:} Bezpłatna 
\end{itemize}

\newpage
\section{TrebEdit} 
\MyFigure[0.3\textwidth]{images/konkurencja/trebedit.png}{Aplikacja TrebEdit}{fig:trebedit}{Google Play}
\begin{itemize}
    
    \item \textbf{Docelowa platforma:} Android
    \item \textbf{Docelowa grupa:} Początkujący web-deweloperzy, osoby uczące się poprzez analizę cudzego kodu 
    \item \textbf{Główne przewagi konkurencyjne:} Możliwość sklonowania dowolnej strony WWW do edycji i eksperymentów, dedykowany pasek skrótów zawierający najpopularniejsze tagi i znaki specjalne, wbudowane mini-kursy i przypomnienia składni dla HTML, CSS, JS oraz PHP, zintegrowana konsola JavaScript i podgląd na żywo.
    \item \textbf{Słabe strony:} Brak wsparcia dla backendu, brak integracji z Git, aplikacja potrafi mieć problemy z asocjacją plików lub ich prawidłowym zapisem. 
    \item \textbf{Model biznesowy:} Bezpłatna z reklamami lub wersja Premium (oferująca m.in. sugestie kodu, sugestie ścieżki plików, emulację urządzenia, niestandardową czcionkę edytora, możliwość edycji paska narzędzi edytora, nielimitowaną ilość otwartych kart, brak reklam): jednorazowa opłata za 169,99 zł lub subskrypcja: miesięczna za 10,99 zł lub roczna za 99,99 zł. 
\end{itemize}
\section{Wnioski}
Z powyższej analizy wynika, że aplikacje naszej konkurencji są często skomplikowane, drogie i nierzadko wykorzystują uciążliwe systemy subskrypcji. Biorąc to pod uwagę, uważamy, że najlepszą decyzją we wczesnej fazie rozwoju projektu będzie skupienie się na głównym module aplikacji - narzędziach do tworzenia stron internetowych. Monetyzacja to również ważny temat - tu ustaliliśmy, że głównym źródłem finansowania będą dobrowolne dotacje od użytkowników. Przewidujemy także jednorazową opłatę za pobranie aplikacji ze sklepu App Store (iOS), jednak będzie to kwota wysoce przystępna.